\paragraph{Proof.}
\emph{Boundary instance.}
Set \(D=2\), \(V=\{1,2\}\), and \(K=1\).  Let
\[
 \eta_{\mathrm{sel}}
 =\bigl(\mathcal G_{\mathrm{sel}},\{S\},(\varnothing,\{1\})\bigr),
 \qquad
 \mathcal G_{\mathrm{sel}}:\ 1\longrightarrow S\longleftarrow2,
\]
where \(S\) is binary and the observed stratum is \(S=1\), and let
\[
 \eta_{\mathrm{dir}}
 =\bigl(\mathcal G_{\mathrm{dir}},\varnothing,(\varnothing,\{1\})\bigr),
 \qquad
 \mathcal G_{\mathrm{dir}}:\ 2\longrightarrow1.
\]
These are the selected-collider and unselected-directed models in the two-variable example of \cite{DaiEtAl2025SelectionIntervention}.  Define the supplied representatives by
\[
 \mathcal M^{(0)}: X_1\mathbin{---}X_2,
 \qquad
 \mathcal M^{(1)}:
 \zeta_1\longrightarrow X_1\longleftarrow X_2,
 \qquad
 \mathbf M=(\mathcal M^{(0)},\mathcal M^{(1)}).
\]
There is no indicator vertex in \(\mathcal M^{(0)}\).

Both candidate graphs are finite and acyclic and have no latent common cause of \(1\) and \(2\).  In \(\eta_{\mathrm{sel}}\), the only selection node is binary, childless, and has only observed parents; in \(\eta_{\mathrm{dir}}\), all selection conditions are vacuous.  Both candidates have observational target \(\varnothing\), nonobservational target \(\{1\}\), and basal pre-treatment selection.  Equip either graph with independent exogenous inputs and a soft intervention that changes only the mechanism of \(X_1\), leaving every other mechanism and exogenous input shared with the basal world.  Thus all member conditions in \(\mathfrak C_D(\mathbf M)\), except the two regime-equivalence matches that we now calculate, hold for both candidates.

For \(\eta_{\mathrm{sel}}\), the indicator-free observational graph has the path
\[
 X_1\longrightarrow S^*\longleftarrow X_2.
\]
After conditioning on \(S^*\), this is an inducing path between the two retained vertices.  Each endpoint is an ancestor of \(S^*\), so the MAG endpoint rule places a tail at each endpoint.  Hence
\[
 \mathcal M_{\eta_{\mathrm{sel}}}^{(0)}:
 X_1\mathbin{---}X_2. \tag{1}
\]
Under target \(\{1\}\), the relevant edges of the interventional twin DAG are
\[
 \zeta_1\to X_1,qquad
 \varepsilon_1\to X_1,qquad
 \varepsilon_1\to X_1^*,qquad
 X_1^*\to S^*,qquad
 X_2\to S^*.
\]
The path
\[
 X_2\to S^*\leftarrow X_1^*\leftarrow\varepsilon_1\to X_1
\]
is inducing after marginalizing \(X_1^*,\varepsilon_1\) and conditioning on \(S^*\).  Because \(X_2\) is an ancestor of \(S^*\) while reality-world \(X_1\) is an ancestor of neither \(X_2\) nor \(S^*\), its projected edge is \(X_2\to X_1\).  The edge \(\zeta_1\to X_1\) is retained.  There is no inducing path from \(\zeta_1\) to \(X_2\): every such path must begin \(\zeta_1\to X_1\leftarrow\varepsilon_1\), where the retained nonendpoint \(X_1\) is a collider having no descendant among the other endpoint or the selected set.  Therefore
\[
 \mathcal M_{\eta_{\mathrm{sel}}}^{(1)}:
 \zeta_1\longrightarrow X_1\longleftarrow X_2. \tag{2}
\]

For \(\eta_{\mathrm{dir}}\), no vertex is marginalized or selected in the observational branch, and consequently
\[
 \mathcal M_{\eta_{\mathrm{dir}}}^{(0)}:
 X_2\longrightarrow X_1. \tag{3}
\]
The target-\(1\) twin graph retains the direct edges \(\zeta_1\to X_1\) and \(X_2\to X_1\).  As there is no selection node, and every path from \(\zeta_1\) to \(X_2\) has the retained collider \(X_1\) as a nonendpoint, the remaining pair is nonadjacent.  Hence
\[
 \mathcal M_{\eta_{\mathrm{dir}}}^{(1)}:
 \zeta_1\longrightarrow X_1\longleftarrow X_2. \tag{4}
\]
The two MAGs in (1) and (3) are Markov equivalent: on two vertices they have the same single adjacency, while there is no triple and hence no unshielded collider or discriminating path.  The MAGs in (2) and (4) are identical.  The two equivalence-match conditions in the definition of \(\mathfrak C_D(\mathbf M)\) therefore hold, proving
\[
 \eta_{\mathrm{sel}},\eta_{\mathrm{dir}}
 \in\mathfrak C_D(\mathbf M). \tag{5}
\]

It remains to verify that the tail at vertex \(2\) is invariant over the \emph{entire} class, rather than merely shared by the two displayed witnesses.  Fix arbitrary \(\eta\in\mathfrak C_D(\mathbf M)\).  Equivalence to \(\mathcal M^{(0)}\) forces the observational adjacency \(X_1-X_2\).  Equivalence to \(\mathcal M^{(1)}\) forces the interventional skeleton
\[
 \zeta_1-X_1-X_2,qquad \zeta_1\not\sim X_2, \tag{6}
\]
and forces the unshielded triple in (6) to be a collider at \(X_1\).

Suppose, for contradiction, that the observational mark at the endpoint of \(X_2\) were an arrowhead.  The observational twin graph has only the two retained system vertices and childless selected sinks, with no latent system vertex.  A non-direct inducing path between \(X_1\) and \(X_2\) must therefore pass through a common selected descendant and gives tails at both endpoints.  Causal sufficiency rules out a bidirected common-cause edge.  Thus an arrowhead at \(X_2\) can arise only from the causal edge
\[
 1\longrightarrow2\quad\text{in }\mathcal G_\eta. \tag{7}
\]
By (6), vertex \(2\) cannot be in \(I_\eta^{(1)}\), since the direct twin edge \(\zeta_1\to X_2\) would make \(\zeta_1\) adjacent to \(X_2\).  The target cannot be empty either, because then the exogenous indicator would be isolated, contrary to the adjacency \(\zeta_1-X_1\) in (6).  Since \(V=\{1,2\}\), it follows that
\[
 I_\eta^{(1)}=\{1\}. \tag{8}
\]
The twin construction copies the edge in (7) into the reality world.  Consequently \(X_1\) is an ancestor of \(X_2\), and the MAG endpoint rule places a tail at \(X_1\) on their interventional edge.  This contradicts the collider at \(X_1\) forced by (6).  The supposition was false.  Since an observational endpoint mark is either a tail or an arrowhead, we have proved
\[
 \operatorname{mark}_{(0,2;12)}
   (\mathcal M_\eta^{(0)})=\mathrm{tail}
 \quad\text{for every }\eta\in\mathfrak C_D(\mathbf M). \tag{9}
\]

At the other endpoint, (1) gives a tail at \(X_1\), whereas (3) gives an arrowhead at \(X_1\).  By the definition of the canonical SI-PAG and (5), that endpoint is noninvariant and hence circular.  By (9), the endpoint at \(X_2\) is invariant and is a tail.  Therefore the canonical observational component has precisely
\[
 \operatorname{mark}_{(0,1;12)}(\mathcal P_{\mathrm{SI}})=\circ,
 \qquad
 \operatorname{mark}_{(0,2;12)}(\mathcal P_{\mathrm{SI}})=\mathrm{tail}. \tag{10}
\]

Finally, soundness of an SI-PAG orientation calculus means that every noncircular mark it outputs is shared by every member of \(\mathfrak C_D(\mathbf M)\).  Orienting the first endpoint as a tail is falsified by \(\eta_{\mathrm{dir}}\), while orienting it as an arrowhead is falsified by \(\eta_{\mathrm{sel}}\).  Hence no sound calculus may orient that endpoint.

This is a finite oracle-graph statement.  No division, inverse, limiting argument, distributional support restriction, rank condition, positivity condition, or overlap condition is used.  The only boundary conditions are \(D=2\), \(K=1\), an indicator-free observational branch, and one nonobservational regime; class nonemptiness is witnessed explicitly by (5). \(\square\)